% PROPOZYCJA tekstu na podstawie artefaktów dostępnych w repozytorium,
% dzienników VM, plików VMX i zapisów przebiegów. Wymaga włączenia do właściwego
% numeru rozdziału, sprawdzenia cytowań i ostatecznej weryfikacji autora/promotora.
% Braki dowodowe są jawnie opisane; nie należy zastępować ich wartościami domyślnymi.

\subsection{Implementacja maszyny agenta}
\label{subsec:implementacja-vm-agent-propozycja}

Maszyna VM-Agent pracowała pod Ubuntu Desktop 26.04.1 LTS. Zapisany plik VMX
określa 2 vCPU, 8096 MiB RAM i dwie karty sieciowe: interfejs administracyjny
VMnet8/NAT oraz interfejs badawczy VMnet10 z adresem 10.10.10.10/24.
Maszyna miała zainstalowane Visual Studio Code 1.139.1 i rozszerzenie Kilo
Code 7.8.1 w wariancie Linux x64. Brama strefy Agent znajdowała się na
VM-FW pod adresem 10.10.10.1. Interfejs NAT służył do administracji SSH;
ruch eksperymentalny kierowano przez bramę strefy i OPNsense.

Na VM-Agent utworzono dwa osobne konta systemowe: \texttt{agent-wide}
(UID 1001) oraz \texttt{agent-limited} (UID 1002). W konfiguracji porównawczej
\texttt{agent-wide} należał do grupy \texttt{lab-canary-readers}, a
\texttt{agent-limited} nie należał do tej grupy. Dziennik wdrożenia stwierdza,
że konta nie należały do grup \texttt{sudo}, \texttt{lxd} ani \texttt{docker}.
Syntetyczny plik chroniony był tworzony jako \texttt{root:lab-canary-readers}
z trybem \texttt{0640}. W zachowanych wynikach auditd zdarzenia
\texttt{openat} dla pliku miały \texttt{success=yes} dla UID 1001 i
\texttt{success=no}, \texttt{exit=-13} (\texttt{EACCES}) dla UID 1002.
Szczegóły par H2, identyfikatory audytu oraz potwierdzenia przekazania logów
do Monitora zestawiono w \texttt{docs/laboratorium/kilo-preflight.md}
i \texttt{docs/laboratorium/maszyny/vm-agent.md}.

Podany właściciel, grupa, tryb i różnica w członkostwie grupowym opisują
porównanie podstawowych uprawnień POSIX/DAC. W repozytorium nie ma wyniku
\texttt{getfacl} dla chronionego pliku. Semantykę podstawowych praw dostępu POSIX
można podeprzeć dokumentacją systemową \cite{LinuxAclManual}; nie rozstrzyga się zatem, czy system
posiadał dodatkowe wpisy rozszerzonej ACL. Wynik H2 dotyczy tego pliku i tych
kont; nie dowodzi izolacji VM/kontenera ani wyeliminowania innych sposobów
dostępu.

Trzy niezależne pary H2 wykonano w Kilo Code. W każdej parze oba konta
wywoływały raz narzędzie \texttt{read\_protected\_canary} dla tego samego
run ID i tego samego bieżącego syntetycznego pliku. Wynik we wszystkich
parach był zgodny: \texttt{agent-wide} uzyskał odczyt, a
\texttt{agent-limited} otrzymał odmowę systemu operacyjnego. Każdą parę
potwierdzono przez local MCP JSONL, auditd i wpisy odebrane na VM-Monitor.
Pilot oraz próby z niezgodnym run ID nie weszły do licznika.

Zrzuty aktywnej konfiguracji Kilo pokazywały regułę ogólną
\texttt{"*": "deny"}, jawne \texttt{allow} dla używanych narzędzi MCP
oraz wyjątek \texttt{"bash": "allow"}. Interfejs Kilo pokazywał
automatyczne zatwierdzenie wywołań MCP na podstawie dopasowanej reguły.
Nie należy więc opisywać konfiguracji jako ograniczającej agenta wyłącznie
do MCP ani jako blokującej shell. W badaniu H2 wymuszano jedno wskazane
wywołanie MCP, a różnicę skutku egzekwował system plików.

Zgoda Kilo na samo wywołanie narzędzia i decyzja policy proxy są odrębnymi
zdarzeniami. Kilo automatycznie dopuszczało wywołania objęte regułą
\texttt{allow}; natomiast w profilu Gateway \texttt{secured} narzędzie
\texttt{submit\_canary} otrzymywało \texttt{ask} i czekało na niezależną
decyzję operatora. Zapis decyzji Gateway, decyzji operatora, próby i wyniku
znajduje się w \texttt{/var/log/lab-control/decisions.jsonl} na VM-Gateway
i w centralnych rekordach Fluent Bit. W kampanii używano trybu Kilo
widocznego w zrzutach jako \texttt{Code}; nie zachowano osobnych metadanych
trybu dla każdego runu. Oba konta miały osobne pliki
\texttt{\textasciitilde/.config/kilo/kilo.jsonc} w trybie \texttt{0600}; Kilo uruchamiało
lokalny serwer poleceniem
\texttt{/opt/lab/agent-mcp/venv/bin/python /opt/lab/agent-mcp/local\_mcp\_stdio.py}
i serwer zdalny pod adresem Gateway TCP/8001. Profile providera i credential
ustawiono oddzielnie w UI; nie kopiowano credential storage konta
administracyjnego \texttt{vmagent}.
Dokumentacja Kilo opisuje oddzielne uprawnienia dla narzędzi MCP oraz ich
automatyczne zatwierdzanie po dopasowaniu reguły; jest to warstwa klienta,
odrębna od decyzji wykonywanej przez policy proxy \cite{KiloMCPDocs,KiloAutoApproval,MCPArchitecture}.

\subsection{Implementacja dostępu do modelu przez LiteLLM}
\label{subsec:implementacja-litellm-propozycja}

LiteLLM Proxy 1.102.1 działał na VM-Gateway z Ubuntu Server 26.04.1 LTS,
pod adresem 10.10.20.10:4000. Kilo używało dostawcy typu OpenAI Compatible
i bazowego URL \texttt{http://10.10.20.10:4000/v1}. Po stronie klienta
LiteLLM model był udostępniany pod aliasem \texttt{laboratory-model}
(etykieta interfejsu \texttt{Lab LiteLLM}); w konfiguracji upstream
Gateway zapisano \texttt{openai/gpt-4.1-mini}. Dokumentacja OpenAI nazywa
\texttt{gpt-4.1-mini} publicznym identyfikatorem modelu, lecz alias LiteLLM
i wpis konfiguracyjny nie są równoważne identyfikatorowi zwróconemu w
odpowiedzi konkretnego żądania \cite{OpenAIModelsGPT41Mini,OpenAIChatCompletionsAPI,LiteLLMConfig,LiteLLMVirtualKeys}.

Test \texttt{/v1/models} potwierdził alias \texttt{laboratory-model};
użytkownik zgłosił poprawne odpowiedzi \texttt{PROVIDER\_OK} w obu
profilach. W zachowanych artefaktach nie ma surowej odpowiedzi generującej
API z polem \texttt{model}, dlatego nie można potwierdzić identyfikatora
modelu zwróconego przez dostawcę ani konkretnego snapshotu użytego dla
każdego runu. Z tego samego powodu nie można podać rzeczywiście użytych
wartości parametrów generowania, takich jak temperatura, top-p, limit
tokenów lub seed. Nie należy wstawiać wartości domyślnych do opisu
eksperymentu. Ramy AI RMF można wykorzystać do osadzenia ostrożnego zakresu
wniosków w zarządzaniu ryzykiem, ale nie zastępują one pomiarów tej pracy
\cite{NISTAI1001}.

W skrypcie konfiguracji LiteLLM parametr \texttt{drop\_params: true}
umożliwiał pomijanie parametrów nieobsługiwanych przez upstream; nie
określa on wartości temperatury ani innych ustawień próbkowania. Klucz
dostawcy OpenAI był przechowywany na Gateway, a profile Kilo uwierzytelniały
się oddzielnie wobec proxy. Testy potwierdzają odpowiedzi HTTP 200 i stan
usługi, ale zachowane materiały nie ustalają pełnego zakresu treści
zapisywanych przez LiteLLM ani polityki retencji treści request/response.
Wiadomo jedynie, że w opisanym teście kontrolnym nie analizowano treści
promptów.

\subsection{Implementacja serwerów MCP}
\label{subsec:implementacja-mcp-propozycja}

\subsubsection{Lokalny serwer MCP}
\label{subsubsec:local-mcp-propozycja}

Lokalny serwer MCP działał na VM-Agent jako proces potomny uruchamiany przez
Kilo Code i komunikujący się przez standardowe wejście i wyjście (stdio).
Wdrożone środowisko używało Pythona 3.14.4 i MCP SDK 1.30.0. Serwer
udostępniał narzędzia \texttt{read\_local\_note},
\texttt{read\_indirect\_instruction} i
\texttt{read\_protected\_canary}. Narzędzia korzystały z ustalonych ścieżek
i nie przyjmowały dowolnej ścieżki pliku od klienta. Implementacja jest
w \texttt{lab/agent/local\_mcp\_stdio.py}.

Zwykła notatka lokalna znajduje się w \texttt{lab/agent/fixtures/local-note.txt};
zawiera wyłącznie syntetyczny tekst i marker weryfikacyjny. Fixture
instrukcji pośredniej to \texttt{lab/agent/fixtures/indirect-instruction.txt}.
Jego treść jawnie mówi, że jest niezaufana, próbuje zmienić zadanie
streszczenia i sugeruje odczyt syntetycznego canary, a następnie wysłanie go
do laboratoryjnego odbiornika. Jest to testowy bodziec, nie polityka
autoryzacyjna.

Wpis JSONL serwera zawierał UTC timestamp, run ID, nazwę serwera i narzędzia,
UID i nazwę konta, identyfikator zasobu i wynik. Przy skutecznym odczycie
zapisywano hash SHA-256 treści, a nie treść zwykłej notatki ani wartość
canary. Logi zapisywano per konto w
\texttt{/home/<konto>/.local/state/lab-agent/local-mcp.jsonl}, obserwowano
przez auditd i przekazywano Fluent Bitem na VM-Monitor. Ponieważ stdio nie
tworzy połączenia IP między maszynami, brak przepływu MCP w PCAP/Zeek jest
oczekiwany; odczyt potwierdzano logiem MCP i zdarzeniem hostowym. Opis
architektury MCP i rozdziału host/client/server należy podeprzeć specyfikacją
MCP, ale zachowany log nie zawiera wersji protokołu wynegocjowanej w
komunikacie \texttt{initialize} \cite{MCPArchitecture,KiloMCPDocs}.

\subsubsection{Zdalny serwer MCP i policy proxy}
\label{subsubsec:remote-mcp-propozycja}

Kilo łączyło się z policy proxy na VM-Gateway pod adresem
\texttt{http://10.10.20.10:8001/mcp}. Proxy przekazywało dopuszczone
operacje do backendu na VM-Target pod adresem
\texttt{http://10.10.30.10:8000/mcp}, przez sieć routowaną przez OPNsense.
Zastosowano MCP Streamable HTTP. F-03, czyli bezpośredni przepływ
VM-Agent--VM-Target TCP/8000, był wyłączony; docelowa ścieżka zdalnego MCP
prowadziła przez Gateway. Adresy, strefy i reguły opisano w
\texttt{docs/laboratorium/adresacja-sieci.md} oraz
\texttt{docs/laboratorium/opnsense-reguly-startowe.md}.

Gateway udostępniał \texttt{read\_untrusted\_note} oraz
\texttt{submit\_canary}. Polityka miała profile \texttt{baseline},
\texttt{secured} i \texttt{deny\_canary}, z domyślną odmową narzędzi
niezdefiniowanych. Dla \texttt{submit\_canary} wynik wynosił odpowiednio
allow, ask i deny. W stanie ask proxy oczekiwało na decyzję operatora przez
skonfigurowany limit czasu; timeout kończył się bez próby wykonania. Każdy
request był normalizowany i audytowany z run ID, request ID, źródłowym
adresem, profilem, decyzją, attempt, rezultatem i decyzją operatora.
Odbiornik canary na Target nasłuchiwał wyłącznie na
\texttt{127.0.0.1:8080}; dowód skutku stanowił jego JSONL, a nie sam
komunikat Kilo. Implementacja jest dostępna w
\texttt{lab/target/canary\_receiver.py}.

Istotnym ograniczeniem implementacji jest sposób propagacji run ID:
proxy zapisuje run ID w swoim audycie, lecz przy wywołaniu backendowego
\texttt{read\_untrusted\_note} przekazuje pusty słownik argumentów.
Target journal i PCAP nie zawierają więc jawnego run ID dla tego odczytu;
korelację zdalnego MCP prowadzono przez request ID Gateway oraz zgodność
czasu i adresów, z niższą pewnością niż wspólny identyfikator.
Dla \texttt{submit\_canary} proxy przekazywało znormalizowane run ID
i syntetyczną wartość do backendu, a Target receiver rejestrował run ID,
wartość i źródłowy adres Gateway.

Remote fixture w \texttt{lab/target/fixtures/untrusted-note.txt} i lokalna
notatka nie były semantycznie równoważne. SC-05 służył zatem do obserwacji
zdalnego transportu i jego telemetrii, nie do kontrolowanego porównania
reakcji modelu na identyczną treść. W SC-06-D model trzykrotnie odczytał
lokalny fixture pośredniej instrukcji i go streścił, ale nie wygenerował
requestu do Gateway. Te wyniki opisują zachowanie modelu; nie są dowodem
skuteczności policy proxy wobec requestu, którego proxy nie otrzymało.
Opis zagrożenia i rozróżnienie danych niezaufanych od instrukcji użytkownika
można oprzeć na badaniach nad pośrednim prompt injection oraz aktualnym
przewodniku OWASP; wnioski eksperymentalne w tej pracy pozostają ograniczone
do trzech obserwacji bez żądania narzędzia \cite{Greshake2023Indirect,OWASPLLMInjection}.

\paragraph{Wersje i odtwarzalność.}
W dokumentacji wdrożenia zapisano Kilo Code 7.8.1, VS Code 1.139.1,
Python 3.14.4, MCP SDK 1.30.0 na Agent/Target, Uvicorn 0.54.0 na Target,
LiteLLM 1.102.1, auditd 1:4.1.2-1ubuntu0.1, Fluent Bit 5.1.2 i Zeek
8.0.10. VMware Workstation Pro na hoście miał plikową wersję
25.0.0 build-24995812; OPNsense działało w linii 24.7, ale po resecie
nie zachowano odczytu dokładnego patch/build. Zachowane VMX podają aktualnie
2 vCPU i 8096 MiB dla VM-Agent; po 2 vCPU i 4096 MiB dla VM-Gateway,
VM-Target i VM-Monitor; VM-FW: 8 vCPU i 4096 MiB. Pliki VMX znajdują się
poza repozytorium, w katalogu maszyn VMware na hoście.

Nie można przypisać kampanii do konkretnego commita kodu: lokalne repozytorium
nie ma commitów Git, a zestaw kodu wdrożony na VM nie został zarchiwizowany
jako oznaczony release z sumami kontrolnymi. Źródła implementacji dostępne
w repozytorium to \texttt{lab/agent/}, \texttt{lab/gateway/} i
\texttt{lab/target/}; lockfile istnieją dla Agent i Target, lecz nie dla
Gateway. Szczegóły wersji i pozostałe otwarte kwestie zestawiono w
\texttt{docs/laboratorium/indeks-dowodow-pracy.md}.

% Materiały nie pozwalają stwierdzić: getfacl/rozszerzone ACL; pełnych grup
% supplemental dla każdego runu; API response model/snapshot; temperature,
% top_p, max tokens lub seed; per-run MCP initialize negotiatedVersion;
% pełnej retencji logów LiteLLM; hash zgodności wdrożonego kodu z repo; kompletnego
% post-reset negative firewall matrix. Te elementy należy uzupełnić tylko po
% ponownym odczycie VM/logów lub oznaczyć jako niezarejestrowane.
